% 論文節の日本語訳（読み物 PDF）用 preamble
% lualatex + jlreq
\documentclass[12pt,paper=a4]{jlreq}
\usepackage{luatexja-fontspec}
\usepackage{amsmath,amssymb,bm}
\usepackage{graphicx}
\usepackage{hyperref}
\usepackage{enumitem}
\usepackage{booktabs}
\usepackage{microtype}
\usepackage{xcolor}
\usepackage{caption}
\usepackage{tcolorbox}
\tcbuselibrary{breakable,skins}

\hypersetup{
  colorlinks=true,
  linkcolor=blue!60!black,
  urlcolor=blue!60!black,
  citecolor=blue!60!black,
}

\graphicspath{{images/}{./}}

\setlength{\parindent}{1em}
\setlength{\parskip}{0.55em}
\setlist[itemize]{leftmargin=1.6em,itemsep=0.25em,topsep=0.35em}
\setlist[enumerate]{leftmargin=1.8em,itemsep=0.25em,topsep=0.35em}

\captionsetup{
  font=small,
  labelfont=bf,
  skip=0.4em,
  justification=raggedright,
  singlelinecheck=false,
}

\newcommand{\vx}{\bm{x}}
\newcommand{\vu}{\bm{u}}
\newcommand{\veps}{\bm{\epsilon}}
\newcommand{\vw}{\bm{w}}

\DeclareMathOperator*{\argmax}{arg\,max}
\DeclareMathOperator*{\argmin}{arg\,min}
\DeclareMathOperator{\Tr}{Tr}
\DeclareMathOperator{\diag}{diag}

% 命題・定理ボックス（ページ跨ぎ時も切れ目を抑える）
\newtcolorbox{propositionbox}[1][]{%
  enhanced,
  breakable,
  colback=blue!4!white,
  colframe=blue!45!black,
  coltitle=white,
  fonttitle=\bfseries\small,
  title={#1},
  attach boxed title to top left={yshift=-2mm, xshift=3mm},
  boxed title style={
    colback=blue!45!black,
    colframe=blue!45!black,
    boxrule=0pt,
    arc=1pt,
  },
  boxrule=0.5pt,
  arc=1.5pt,
  left=5pt, right=5pt, top=8pt, bottom=5pt,
  before skip=0.9em,
  after skip=0.9em,
  pad at break=3mm,
}

% 図はフロートせず本文位置に置く（分断を防ぐ）
\newcommand{\translationfigure}[2]{%
  \begin{center}
    \includegraphics[width=0.88\linewidth]{#1}\\[0.35em]
    {\small #2\par}
  \end{center}
  \vspace{0.4em}%
}

\begin{document}
{\color{gray}\footnotesize 出典: §5.4\par}
\vspace{0.2em}
\begin{propositionbox}[原題]
\textbf{原題}: (Optional) Properties of the Canonical Affine Flow
\textbf{日本語タイトル}: （任意）正準アフィンフローの性質
\textbf{出典}: §5.4 of *The Principles of Diffusion Models*
\end{propositionbox}


\section*{5.4 （任意）正準アフィンフローの性質}


2つの端点分布 $p_0 = p_{\mathrm{src}}$ と $p_1 = p_{\mathrm{tgt}}$ が与えられたとき、フローマッチング（FM）（Lipman et al., 2022）および整流フロー（Rectified Flow, RF）（Liu, Gong, et al., 2022）における条件付き経路を定義する、自然で広く用いられる選択は線形補間

\[
a_t = 1 - t, \qquad b_t = t
\]


であり、これにより補間

\[
x_t = (1-t)x_0 + t x_1, \qquad x_0 \sim p_{\mathrm{src}},\ x_1 \sim p_{\mathrm{tgt}}
\]


が得られる。この選択の下で学習目的は次に簡略化される：

\[
\mathbb{E}_{t\sim\mathcal{U}[0,1]}\,\mathbb{E}_{x_0,x_1}\Bigl[\bigl\|v_\phi(x_t, t) - (x_1 - x_0)\bigr\|_2^2\Bigr].
\]


この線形フローはいくつかの魅力的な性質をもつ。特に、Reflowとして知られる反復的洗練スキームを許容し、周辺を保ちつつ分布間の経路を次第に直線化する。

\subsection*{5.4.1 フローの整流：ノイズの多い推測から構造化された対付けへ}


\textbf{なぜカップリングを変えるのか？} 独立カップリング

\[
\pi(x_0, x_1) = p_{\mathrm{src}}(x_0)\, p_{\mathrm{tgt}}(x_1)
\]


の下では、各条件付き補間 $(1-t)x_0 + t x_1$ は直線である。しかし §5.4 で強調したように、対応する周辺速度は無関係な端点対からの方向を平均するため、誘導される周辺ODE軌道は直線である必要がない。そのような曲がった周辺ダイナミクスは、少数の数値ステップだけで精度よく積分するのが難しくなり得る。
整流（rectification）はこの問題に、端点カップリングを変えることで対処する。学習済みODEが各ソースサンプルをモデル生成の端点へ輸送し、依存カップリングを生じ、それが新しい条件付き補間の定義に用いられる。

\textbf{依存カップリングによるフローの整流.} 任意の対付けに頼るのではなく、事前学習された拡散モデル $v_{\phi^\times}(\cdot, t)$ をPF-ODEのドリフトとして用い、各ソース点を決定論的に輸送する。$z(0) = z_0 \sim p_{\mathrm{src}}$ から出発し、

\[
\frac{dz(t)}{dt} = v_{\phi^\times}(z(t), t), \qquad t \in [0, 1]
\]


を積分して、事前学習モデルが学習したデータ空間近傍に位置する $\hat{z}_1 := z(1)$ を得る。結果として得られる対 $(z_0, \hat{z}_1)$ は依存カップリングをなす：任意の補間ではなく、構造化されたモデル誘導経路に従う。この考え方は、自然にアフィン参照経路 $x_t = a_t x_0 + b_t x_1$（$x_0 \sim p_{\mathrm{src}}$、$x_1 \sim p_{\mathrm{tgt}}$）へ拡張される。

\textbf{アルゴリズム 2: Rectify操作}

\textbf{入力:} 参照経路 $\{x_t\}_{t\in[0,1]}$（例：$x_t = a_t x_0 + b_t x_1$）

\begin{enumerate}
\item \textbf{拡散の事前学習.} 選んだ経路上で次を最小化して $v_{\phi^\times}$ を適合させる：
\end{enumerate}

\[
\phi^\times \in \arg\min_\phi\, \mathbb{E}_{t,x_0,x_1}\Bigl[\Bigl\|v_\phi(x_t, t) - \frac{dx_t}{dt}\Bigr\|_2^2\Bigr].
\]


\begin{enumerate}
\item \textbf{整流.} $z_0 \sim p_{\mathrm{src}}$ をサンプルし、
\end{enumerate}

\[
\frac{dz(t)}{dt} = v_{\phi^\times}(z(t), t), \qquad z(0) = z_0, \qquad t \in [0, 1]
\]

を積分して $\hat{z}_1 = z(1)$ と軌道 $\{z(t)\}_{t\in[0,1]}$ を得る。

\textbf{出力:} 依存的（一貫した）対 $(z_0, \hat{z}_1)$、または完全な軌道。

\textbf{なぜうまくいくか：周辺保存構造.} $\Phi_{0\to t}$ を、事前学習された拡散 $v_{\phi^\times}$ が定める上記ODEが生成するフロー写像とする。すると $z(t) = \Phi_{0\to t}(z_0)$、$\hat{z}_1 = \Phi_{0\to 1}(z_0)$ である。Rectify手続きは各ソース点をそのフロー端点と対付け、決定論的同時分布

\[
\pi_{\mathrm{Rectify}}(z_0, z_1) = p_{\mathrm{src}}(z_0)\, \delta\bigl(z_1 - \Phi_{0\to 1}(z_0)\bigr)
\]


を与える。直ちに次の2つの帰結がある：

\begin{itemize}
\item \textbf{ソース周辺は保存される:} $\displaystyle\int \pi_{\mathrm{Rectify}}(z_0, z_1)\, dz_1 = p_{\mathrm{src}}(z_0)$。
\item \textbf{フローに沿ったプッシュフォワード:} $(\Phi_{0\to t})_{\#} p_{\mathrm{src}} = \mathrm{Law}(z(t))$、すなわち時刻 $t$ の分布は $p_{\mathrm{src}}$ の $\Phi_{0\to t}$ によるプッシュフォワードである。
\end{itemize}


$v_{\phi^\times}$ が与えられた参照経路 $x_t$ のオラクルドリフトに一致するなら、すべての中間周辺が一致する：

\[
\mathrm{Law}(z(t)) = \mathrm{Law}(x_t), \quad \forall t\in[0,1], \qquad\text{かつ}\qquad (\Phi_{0\to 1})_{\#} p_{\mathrm{src}} = p_{\mathrm{tgt}}.
\]


\textbf{要約.} 整流は独立な端点カップリングを、学習されたフローが誘導するカップリングで置き換える。反復的なReflowの下で、これは少数の数値ステップで積分しやすい周辺輸送を次第に生じさせ得る。
正準経路に対し、Rectifyを繰り返し適用する（「Reflow」）と、輸送コストを増やさずに軌道をさらに直線化し、学習をさらに容易にする。

\subsection*{5.4.2 Reflow：フローの反復的直線化}


\textbf{なぜReflowか？} 独立な対付けはしばしば $p_{\mathrm{src}}$ と $p_{\mathrm{tgt}}$ の間に不規則で蛇行するODE軌道を誘導し、シミュレーション中の離散化誤差と分散を増大させる。これは自然な問いを提起する：

\begin{propositionbox}[問い 5.4.1]
正しい周辺を保ちつつ、2分布間でより直線に近い輸送経路を誘導するカップリングを学習できるか？
\end{propositionbox}


これがReflowを動機づける：Rectifyを繰り返し適用してカップリングを更新し、後続のフローが積分しやすくなるようにする。

\textbf{核心的アイデア：Rectifyによる再帰的直線化.} 積カップリング $\pi^{(0)} := p_{\mathrm{src}}(x_0)p_{\mathrm{tgt}}(x_1)$ 上の正準補間

\[
x_t = (1-t)x_0 + t x_1
\]


から出発する。Rectifyを適用すると独立対付けが、学習済みODEが誘導する依存対 $(z_0, \hat{z}_1)$ で置き換えられる。この手続きを繰り返すこと——Reflow——は、少数の数値ステップで精度よく積分しやすい、より直線的な軌道を生じる傾向がある。厳密な直線化保証は後の §5.4.3 で議論する。

\textbf{Reflow手続き.} 各反復は2ステップを行う：

\begin{itemize}
\item \textbf{フローの再適合:} 現在のカップリングのサンプルから新しい速度場を学習する：
\end{itemize}

\[
\phi_{k+1} = \arg\min_\phi\, \mathcal{L}\bigl(\phi; \pi^{(k)}\bigr),
\]

ここで
\[
\mathcal{L}\bigl(\phi; \pi^{(k)}\bigr) := \mathbb{E}_{t,(z_0^{(k)},\hat{z}_1^{(k)})\sim\pi^{(k)}}\Bigl[\bigl\|v_\phi(z_t, t) - (\hat{z}_1^{(k)} - z_0^{(k)})\bigr\|_2^2\Bigr], \quad (5.4.1)
\]

$z_t = (1-t)z_0^{(k)} + t\hat{z}_1^{(k)}$。

\begin{itemize}
\item \textbf{新しいカップリングの生成:} 新しいソースサンプル $z_0^{(k+1)} \sim p_{\mathrm{src}}$ から出発して学習済みODEを解く：
\end{itemize}

\[
\hat{z}_1^{(k+1)} \leftarrow z_0^{(k+1)} + \int_0^1 v_{\phi_{k+1}}(z(t), t)\, dt,
\]

更新されたカップリングを
\[
\pi^{(k+1)}(z_0, z_1) := p_{\mathrm{src}}(z_0)\, \delta\bigl(z_1 - \hat{z}_1^{(k+1)}\bigr)
\]

と定義する。

言い換えれば、ReflowはRectify演算子の繰り返し適用とみなせ、次第に洗練されたカップリングの列を生じる：

\[
\pi^{(k+1)} = \mathrm{Rectify}\bigl(\pi^{(k)}\bigr). \quad (5.4.2)
\]


これによりフローとカップリングがともに発展し、次第に安定した輸送経路が得られる。

\subsection*{5.4.3 Reflowの性質}


Reflowの有用性を駆動する2つの鍵となる理論的性質がある：輸送コストを減らし、軌道を直線化することである。

\subsubsection*{I. Reflowは輸送コストを決して増やさない}


$c(y)$ を凸コスト関数とする（例：$p \ge 1$ に対する $\|y\|_2^p$）。各Rectifyステップは、コストが元のものより悪くない新しいカップリング $(z_0, \hat{z}_1)$ を形成する：

\begin{propositionbox}[命題 5.4.1（Rectifyは輸送コストを減らし得る）]
理想的な速度場 $v^* = v_{\phi^\times}$ を仮定すると、
\[
\mathbb{E}\bigl[c(\hat{z}_1 - z_0)\bigr] \le \mathbb{E}\bigl[c(x_1 - x_0)\bigr]
\]
が成り立つ。
\end{propositionbox}


\textbf{命題の証明.}
Jensenの不等式から従う。完全な導出は Liu, Gong, et al.（2022）を参照。∎

この結果を再帰的に適用すると、Reflow過程が輸送コストを増やさないことがわかる。

\subsubsection*{II. Reflowは経路を直線化する}


Reflowの再帰的適用はますます直線的な軌道を促進するが、直線性の尺度が各反復で単調に改善することは保証されない。これを測るため、経路 $Y = \{y_t\}_{t\in[0,1]}$ の直線性汎関数を

\[
S(Y) := \int_0^1 \mathbb{E}\Bigl[\Bigl\|y_1 - y_0 - \frac{dy_t}{dt}\Bigr\|_2^2\Bigr]\, dt
\]


と定義する。$S(Y) = 0$ ならば $Y$ は厳密に直線である。

\begin{propositionbox}[命題 5.4.2（Reflowは確率的経路を直線化する）]
整流された経路 $Z^{(k)}$ に対し、
\[
\min_{k\in\{0,\ldots,K\}} S(Z^{(k)}) \le \frac{\mathbb{E}\bigl[\|x_1 - x_0\|^2\bigr]}{K}
\]
が成り立つ。
\end{propositionbox}


\textbf{命題の証明.}
Liu, Gong, et al.（2022）の定理 3.7 を参照。∎

\translationfigure{fig-5-11.jpg}{図5.11: Reflowの模式図。Rectify手続きにより経路は次第に直線的になる。}


\emph{図5.11: Reflowの模式図。Rectify手続きにより経路は次第に直線的になる。}
出典: Adapted from Liu, Gong, et al. (2022).

線形補間カーネルを用いるFMまたはRFの定式化は、Reflow手続きと合わせて、より単純な学習目的と、確率的カップリングを洗練する実用的方法を提供する。理論的詳細については Liu, Gong, et al.（2022）; Liu（2022）を参照されたい。

\subsubsection*{III. 最適輸送との接続}


最後に、直線的カップリングが最適輸送（Optimal Transport, OT）の意味で必ずしも最適ではないことに注意する。これは §7.2 で導入される用語を含むため、OTに不慣れな読者はその節を参照されたい。
二次コストの最適輸送の特徴は、粒子が直線に沿って進むことである：$x_0$ にある粒子は、$T$ を最適輸送写像として $x_t = (1-t)x_0 + t\, T(x_0)$ により $T(x_0)$ へ動く。しかし、そのような直線経路を生成する任意の写像 $S$、すなわち $x_t = (1-t)x_0 + t\, S(x_0)$ が最適であるとは限らない。写像 $S$ が最適フローを与えるのは、Mongeコスト $\mathbb{E}[\|x_0 - S(x_0)\|^2]$ を最小化するときだけである。したがって直線経路は必要だが十分ではなく、最適性は正しい端点写像 $T$ にも依存する。

\begin{propositionbox}[例：直線カップリングは最適である必要がない]
$p_{\mathrm{src}} = p_{\mathrm{tgt}} = \mathcal{N}(0, I)$ とする。コスト $c(x, y) = \|x - y\|^p$（$p > 0$）に対し、$c$-最適カップリングは恒等カップリング $\pi^*$ であり、$\pi^*$ は $x \sim p_{\mathrm{src}}$ に対する $(x, x)$ の法則である。
次に、カップリング $\pi_A$ を $(x, Ax)$ の法則として考える。ここで $x \sim p_{\mathrm{src}}$、$A$ は $A^\top A = I$、$\det(A) = 1$、$A \ne I$、かつ $-1$ が固有値でない回転行列とする。すると $\pi_A$ は $p_{\mathrm{src}}$ と $p_{\mathrm{tgt}}$ の有効なカップリングであり、$x$ と $Ax$ の間の直線経路に対応するが、可逆ヘッセをもつ任意の二回微分可能な狭義凸コスト $c$ に対して $c$-最適ではない。劣最適性は回転変換から生じる。式 (7.5.2) で議論されるように、回転を除いても最適カップリングに至るとは限らない。∎
\end{propositionbox}


OTとの接続は §7.5.2 で引き続き探る。
\end{document}
